\section{Análisis de requerimientos}
\label{sec:requerimientos}

\subsection{Metodología del levantamiento y fuentes}

Cada requerimiento declara su origen. Se emplean cuatro tipos de fuente, con distinto peso probatorio.

\begin{longtable}{@{}C{1.5cm}L{5.2cm}L{7.0cm}@{}}
\filaencab \encab{Código} & \encab{Fuente} & \encab{Peso probatorio} \\ \midrule \endhead
\textbf{[PO]} & Solicitud del propietario del proyecto & Alto: es requerimiento declarado, no inferido \\
\textbf{[EM]} & Estudio de mercado: oferta regional y sector de alarmas monitoreadas & Medio-alto: evidencia de lo que el mercado cubre y no cubre \\
\textbf{[BT]} & Base teórica: estándares SIA DC-05 y DC-09, LoRaWAN RP002, ISO/IEC 25010, RFC 8428 & Alto para requerimientos técnicos \\
\textbf{[DP]} & Derivación del problema: consecuencia lógica de una causa identificada & Medio: depende de que la causa esté bien fundada \\
\bottomrule
\caption{Tipos de fuente empleados en el levantamiento.}
\label{tab:fuentes}
\end{longtable}

\begin{cajahallazgo}[title={Brecha de fuente declarada}]
Falta la fuente \textbf{[ENT]}: entrevistas directas a las partes interesadas. Hasta obtenerlas, los
requerimientos marcados \textbf{[PO]} se sustentan en la indicación transmitida por el profesor guía y
no en levantamiento directo. Esta limitación se declara explícitamente, y su cierre constituye la acción
prioritaria del período.
\end{cajahallazgo}

\subsection{Estructura modular de la solución}

El análisis descompone la solución en ocho módulos con responsabilidades separadas. Esta descomposición
condiciona posteriormente la elección de la tecnología de implementación (sección~\ref{sec:lenguaje}) y
el diseño de la arquitectura de la Etapa 2.

\begin{longtable}{@{}C{1.4cm}L{5.4cm}C{1.6cm}L{4.6cm}@{}}
\filaencab \encab{Módulo} & \encab{Responsabilidad} & \encab{Req.} & \encab{Causas que aborda} \\ \midrule \endhead
\textbf{M1} & Ingesta multiprotocolo & 10 & C1, C2 \\
\textbf{M2} & Normalización de eventos y telemetría & 7 & C1, C5 \\
\textbf{M3} & Inventario y aprovisionamiento de dispositivos & 6 & C4, C5 \\
\textbf{M4} & Supervisión de salud de enlace & 7 & \textbf{C5} \\
\textbf{M5} & Persistencia y consulta & 5 & C1, C3, C4 \\
\textbf{M6} & Alertas y notificación & 4 & C5 \\
\textbf{M7} & Seguridad, identidad y multi-tenencia & 7 & C1, C3, C4 \\
\textbf{M8} & Interfaz de administración & 6 & C1, C2, C4, C5 \\
\bottomrule
\caption{Módulos de la solución y su relación con las causas del problema.}
\label{tab:modulos}
\end{longtable}

\subsection{Síntesis del catálogo}

El catálogo completo, con el enunciado, la fuente, la causa y la prioridad de cada elemento, se presenta
en el anexo~\ref{anx:requerimientos}. Se cataloga un total de \textbf{68 requerimientos}, frente al
mínimo de 20 comprometido en OE1.

\begin{longtable}{@{}L{5.4cm}C{2.4cm}C{2.6cm}C{2.2cm}@{}}
\filaencab \encab{Prioridad} & \encab{Funcionales} & \encab{No funcionales} & \encab{Total} \\ \midrule \endhead
\textbf{Must}: indispensable para el producto & 35 & 14 & \textbf{49} \\
\textbf{Should}: importante, postergable & 14 & 2 & 16 \\
\textbf{Could}: deseable & 3 & 0 & 3 \\
\midrule
\textbf{Total catalogado} & \textbf{52} & \textbf{16} & \textbf{68} \\
\bottomrule
\caption{Distribución de requerimientos según priorización MoSCoW.}
\label{tab:moscow}
\end{longtable}

Los requerimientos no funcionales se clasifican según las características de calidad de la norma
\textbf{ISO/IEC 25010}: eficiencia de desempeño, fiabilidad, seguridad, mantenibilidad, y portabilidad
y compatibilidad.

\subsection{Requerimientos que anclan el trabajo al enfoque de redes}

Cuatro requerimientos no funcionales merecen mención destacada por vincular el trabajo con el dominio de
administración de infraestructura de red.

\begin{longtable}{@{}L{1.8cm}L{11.9cm}@{}}
\filaencab \encab{ID} & \encab{Requerimiento y justificación} \\ \midrule \endhead
\textbf{RNF-09} & Los dispositivos IoT deben poder residir en un \textbf{segmento de red separado} del
segmento de usuarios, con acceso restringido a los puertos de ingesta. Corresponde a segmentación de red
e impide que un sensor comprometido alcance la red corporativa. \\
\textbf{RNF-10} & El sistema \textbf{no debe requerir puertos entrantes abiertos hacia los
dispositivos}: la conexión se inicia siempre desde el dispositivo hacia la plataforma. Resuelve la
alcanzabilidad de dispositivos situados tras traducción de direcciones y cortafuegos, sin exponerlos. \\
\textbf{RNF-12} & Añadir un transporte nuevo debe requerir implementar solo un adaptador, sin modificar
el núcleo de normalización. Es la traducción arquitectónica de la indicación de no cerrarse a una sola
tecnología: si añadir un transporte obliga a modificar el núcleo, el sistema no es multiprotocolo sino
monoprotocolo con excepciones. \\
\textbf{RNF-16} & El mensaje de evento debe caber en el payload de LoRaWAN AU915 en su configuración más
restrictiva, de 11 bytes. Verificado empíricamente: 8 bytes (sección~\ref{sec:tamano}). \\
\bottomrule
\caption{Requerimientos no funcionales de anclaje al dominio de redes.}
\label{tab:rnfredes}
\end{longtable}

\subsection{Requerimientos que sustentan la diferenciación}

Del análisis de la oferta regional se desprende que el valor diferencial no se distribuye
uniformemente. Los siguientes requerimientos corresponden a capacidades que ningún proveedor regional
identificado satisface.

\begin{longtable}{@{}L{2.9cm}L{10.8cm}@{}}
\filaencab \encab{Requerimiento} & \encab{Fundamento de la diferenciación} \\ \midrule \endhead
\textbf{RF-04} Receptor DC-09 nativo & Permite ingerir paneles de alarma ya instalados sin reemplazarlos \\
\textbf{RF-24, 27, 30} Salud de enlace & La oferta regional presenta el dato del sensor, no el estado del dispositivo ni de su enlace \\
\textbf{RF-34} Exportación abierta & Materializa la ausencia de dependencia de proveedor: una plataforma que se declara abierta y no permite retirar los datos reproduce el problema que pretende resolver \\
\textbf{RF-51} Vista por transporte & Ninguna plataforma presenta la flota organizada por tecnología de red \\
\textbf{RNF-12} Extensibilidad & Las soluciones cerradas no admiten hardware de terceros \\
\textbf{RNF-15} Solo componentes abiertos & Elimina el costo de licencia por dispositivo \\
\bottomrule
\caption{Requerimientos que sustentan la diferenciación competitiva.}
\label{tab:diferenciacion}
\end{longtable}

\subsection{Verificación de cobertura}

\begin{longtable}{@{}L{6.4cm}C{2.2cm}L{4.6cm}@{}}
\filaencab \encab{Causa del problema} & \encab{Req.} & \encab{Módulos principales} \\ \midrule \endhead
C1 Fragmentación de protocolos y formatos & 24 & M1, M2, M7 \\
C2 Acoplamiento hardware-plataforma & 5 & M1, M8 \\
C3 Costo y dependencia comercial & 5 & M5, M7 \\
C4 Ausencia de capacidad técnica interna & 11 & M3, M8 \\
C5 Falta de visibilidad de la conectividad & 20 & M4, M6 \\
\bottomrule
\caption{Cobertura de las causas del problema por el catálogo de requerimientos.}
\label{tab:coberturareq}
\end{longtable}

Ninguna causa queda sin requerimientos asociados, y ningún requerimiento del catálogo carece de causa
declarada.

\subsection{Evaluación crítica de alternativas de solución}

La evaluación de la decisión de mayor peso técnico ---el protocolo de comunicación--- se desarrolla en la
sección~\ref{sec:protocolo}. Esta subsección cubre las decisiones restantes.

\subsubsection{Servidor de red LoRaWAN}

\begin{longtable}{@{}L{3.0cm}L{4.6cm}L{4.4cm}C{1.8cm}@{}}
\filaencab \encab{Alternativa} & \encab{Ventajas} & \encab{Desventajas} & \encab{Decisión} \\ \midrule \endhead
\textbf{ChirpStack} & Abierto, autoalojable, publica en MQTT de forma nativa, gestión completa de dispositivos & Requiere operar el componente & \textbf{Seleccionado} \\
The Things Stack & Funcionalmente equivalente, comunidad amplia & Despliegue autoalojado más exigente en recursos & Descartado \\
The Things Network (servicio público) & Sin operación propia, gratuito & Los datos transitan por infraestructura de terceros: contradice RNF-14 y el control del dato & Descartado \\
Desarrollo propio & Control total & Reimplementar LoRaWAN no es el aporte del trabajo y consumiría todo el presupuesto de tiempo & Descartado \\
\bottomrule
\caption{Evaluación de alternativas de servidor de red LoRaWAN.}
\label{tab:lorawan}
\end{longtable}

\subsubsection{Motor de persistencia}

\begin{longtable}{@{}L{3.0cm}L{4.6cm}L{4.4cm}C{1.8cm}@{}}
\filaencab \encab{Alternativa} & \encab{Ventajas} & \encab{Desventajas} & \encab{Decisión} \\ \midrule \endhead
\textbf{TimescaleDB} (extensión de PostgreSQL) & Un solo motor para datos relacionales y series de tiempo; SQL estándar; licencia abierta & Menos optimizado que InfluxDB en escritura masiva pura & \textbf{Seleccionado} \\
InfluxDB & Muy eficiente en series de tiempo & Obliga a operar dos motores, porque el inventario y la identidad son relacionales & Descartado \\
PostgreSQL sin extensión & Máxima simplicidad & Sin optimización de series de tiempo ni políticas de retención automáticas & Descartado \\
\bottomrule
\caption{Evaluación de alternativas de motor de persistencia.}
\label{tab:bd}
\end{longtable}

\begin{cajaclave}[title={Razón de la selección}]
El factor decisivo es operativo y no de rendimiento. Se trata de un proyecto unipersonal con
aproximadamente 15 horas semanales: operar, respaldar y consultar dos motores de base de datos distintos
duplica la carga de mantenimiento sin beneficio proporcional a la escala del prototipo. El volumen
esperado ---decenas de dispositivos reportando cada pocos minutos--- se sitúa muy por debajo del punto
en que la ventaja de InfluxDB resulta relevante. \textbf{Se documenta que a escala mayor la decisión
podría invertirse}, lo que constituye una recomendación de trabajo futuro.
\end{cajaclave}

\subsubsection{Broker MQTT y estrategia de despliegue}

Se selecciona \textbf{Mosquitto} frente a EMQX por proporcionalidad: sus capacidades de agrupamiento y
alta concurrencia resultan desproporcionadas para la escala del prototipo, mientras que Mosquitto
provee cifrado del transporte y autenticación con consumo mínimo.

Se selecciona el despliegue mediante \textbf{contenedores orquestados}, frente a la instalación directa
en el sistema operativo. La razón no es preferencia técnica sino requisito de evaluación: la Etapa 4
exige demostrar el funcionamiento \enquote{bajo condiciones operativas realistas}, y que el sistema
completo se levante con un solo comando en una máquina distinta a la de desarrollo constituye evidencia
directa de ello.

\subsubsection{Lenguaje y marco de trabajo de la aplicación}
\label{sec:lenguaje}

\begin{longtable}{@{}C{1.4cm}L{7.4cm}C{1.6cm}@{}}
\filaencab \encab{} & \encab{Criterio} & \encab{Peso} \\ \midrule \endhead
\textbf{L-A} & Experiencia previa del desarrollador & 25\% \\
\textbf{L-B} & Correspondencia entre estructura del código y arquitectura documentable & 15\% \\
\textbf{L-C} & Generación automática de la especificación de la interfaz & 15\% \\
\textbf{L-D} & Unificación del lenguaje con la interfaz de usuario & 15\% \\
\textbf{L-E} & Madurez de bibliotecas para MQTT, PostgreSQL y manipulación binaria & 15\% \\
\textbf{L-F} & Tipado estático & 10\% \\
\textbf{L-G} & Rendimiento en tiempo de ejecución & 5\% \\
\bottomrule
\caption{Criterios de evaluación de la tecnología de implementación y su ponderación.}
\label{tab:criterioslenguaje}
\end{longtable}

\begin{longtable}{@{}L{4.0cm}C{1.1cm}C{1.1cm}C{1.1cm}C{1.1cm}C{1.1cm}C{1.1cm}C{1.1cm}C{1.4cm}@{}}
\filaencab \encab{Alternativa} & \encab{L-A} & \encab{L-B} & \encab{L-C} & \encab{L-D} & \encab{L-E} & \encab{L-F} & \encab{L-G} & \encab{Total} \\ \midrule \endhead
\textbf{TypeScript + NestJS} & 4 & 5 & 5 & 5 & 5 & 5 & 4 & \textbf{4,70} \\
TypeScript + Fastify & 5 & 3 & 4 & 5 & 5 & 5 & 5 & 4,55 \\
Python + FastAPI & 3 & 3 & 5 & 2 & 5 & 4 & 3 & 3,55 \\
Java + Spring Boot & 2 & 5 & 4 & 1 & 5 & 5 & 4 & 3,45 \\
Go + Echo & 1 & 3 & 3 & 1 & 4 & 5 & 5 & 2,65 \\
\bottomrule
\caption{Matriz comparativa ponderada de tecnologías de implementación.}
\label{tab:matrizlenguaje}
\end{longtable}

\textbf{Sobre el criterio L-B, que decide el resultado.} El análisis identificó ocho módulos con
responsabilidades separadas (tabla~\ref{tab:modulos}). NestJS impone una estructura de módulos con
dependencias explícitas que corresponde uno a uno con esa descomposición: cada módulo del análisis se
convierte en un módulo del código, de modo que los diagramas de arquitectura de la Etapa 2 describen la
estructura real del sistema. Es la diferencia entre documentar lo que existe y confiar en que lo que
existe se asemeje a lo documentado.

\textbf{Sobre el criterio L-D.} El evento Contact ID posee una estructura estricta de siete campos. Con
un único lenguaje en todo el sistema, esa estructura se define una sola vez y la emplean el validador de
ingesta, la capa de persistencia, la interfaz de programación y la interfaz de usuario. Con dos
lenguajes debe definirse dos veces y mantenerse sincronizada manualmente, lo que constituye una fuente
permanente de error.

\textbf{Sobre el bajo peso de L-G.} El requerimiento RNF-01 exige sostener 50 eventos por segundo, cifra
que todas las alternativas evaluadas superan holgadamente. Ponderar el rendimiento por encima de la
velocidad de desarrollo equivaldría a optimizar el recurso abundante a costa del escaso, que en este
proyecto es el tiempo del desarrollador.

\begin{cajaclave}[title={Decisión}]
Se selecciona \textbf{TypeScript con NestJS}, con la interfaz de administración en el mismo lenguaje y
repositorio, compartiendo las definiciones de tipos del dominio.

Se declara que la diferencia con Fastify, de 0,15 puntos, \textbf{no constituye una ventaja decisiva}:
ambas opciones son defendibles. Si durante la Etapa 2 la curva de aprendizaje de NestJS resulta más
costosa de lo previsto, migrar a Fastify constituiría una desviación aceptable y quedaría documentada
como tal.
\end{cajaclave}

\subsection{Restricciones y supuestos}

\begin{longtable}{@{}L{1.4cm}L{5.6cm}L{6.6cm}@{}}
\filaencab \encab{} & \encab{Restricción} & \encab{Consecuencia} \\ \midrule \endhead
\textbf{RE-01} & Presupuesto de aproximadamente 15 horas semanales & El alcance se limita al núcleo demostrativo \\
\textbf{RE-02} & Proyecto unipersonal & Se descartan alternativas que exijan operar múltiples componentes complejos \\
\textbf{RE-03} & Sin licencias de pago & Solo componentes de código abierto \\
\textbf{RE-04} & Banda 915--928\,MHz, plan AU915, sin licencia & Payload limitado; verificado como suficiente \\
\textbf{RE-05} & Umbral de aprobación del 70\% de logro & El alcance se dimensiona para superar el umbral, no para maximizar funcionalidad \\
\textbf{RE-06} & Disponibilidad de hardware no confirmada & La validación debe funcionar con emulación por software \\
\bottomrule
\caption{Restricciones del proyecto.}
\label{tab:restricciones}
\end{longtable}
